
\documentclass[12pt,a4paper]{article}

% =========================================================
% INDIAN PATENT APPLICATION - COMPLETE SPECIFICATION
% Drafting template based on Form 2 and Section 10
% of the Patents Act, 1970, and Rule 13 of the Patents
% Rules, 2003. Verify current official requirements.
% =========================================================

\usepackage[a4paper,margin=25mm]{geometry}
\usepackage{setspace}
\usepackage{graphicx}
\usepackage{amsmath,amssymb}
\usepackage{enumitem}
\usepackage{array}
\usepackage{longtable}
\usepackage{fancyhdr}
\usepackage{caption}
\usepackage{float}
\usepackage{hyperref}

\setstretch{1.25}
\setlength{\parindent}{0pt}
\setlength{\parskip}{6pt}
\setcounter{secnumdepth}{0}

% Fill in the following details
\newcommand{\InventionTitle}
{[TITLE OF THE INVENTION]}
\newcommand{\ApplicantName}
{[FULL LEGAL NAME OF APPLICANT]}
\newcommand{\ApplicantNationality}
{[NATIONALITY]}
\newcommand{\ApplicantAddress}
{[FULL POSTAL ADDRESS WITH PIN CODE, STATE, COUNTRY]}
\newcommand{\InventorNames}
{[FULL NAMES OF ALL INVENTORS]}
\newcommand{\FilingDate}
{[DATE, IF APPLICABLE]}
\newcommand{\ApplicationNumber}
{[APPLICATION NUMBER, IF ALREADY ALLOTTED]}

\begin{document}

% =========================================================
% FORM 2: FRONT PAGE
% =========================================================

\begin{center}
    \textbf{\Large FORM 2}\\[6pt]
    \textbf{THE PATENTS ACT, 1970}\\
    \textbf{(39 of 1970)}\\[4pt]
    \textbf{AND}\\[4pt]
    \textbf{THE PATENTS RULES, 2003}\\[8pt]
    \textbf{COMPLETE SPECIFICATION}\\[4pt]
    (See section 10 and rule 13)
\end{center}

\vspace{8mm}

\textbf{1. TITLE OF THE INVENTION}

\InventionTitle

\vspace{4mm}

\textbf{2. APPLICANT(S)}

\begin{enumerate}[label=(\alph*)]
    \item \textbf{Name:} \ApplicantName
    \item \textbf{Nationality:} \ApplicantNationality
    \item \textbf{Address:} \ApplicantAddress
\end{enumerate}

% Add additional applicant particulars if required.

\vspace{4mm}

\textbf{3. PREAMBLE TO THE DESCRIPTION}

The following specification particularly describes the
invention and the manner in which it is to be performed.

\vspace{5mm}

\textbf{4. DESCRIPTION}

The description commences on the following page.

\vspace{3mm}

\textbf{5. CLAIMS}

The claims defining the scope of the invention are set out
in the claims section of this complete specification.

\vspace{3mm}

\textbf{6. DATE AND SIGNATURE}

Date: \FilingDate

Signature: \rule{55mm}{0.4pt}

Name: [APPLICANT OR AUTHORISED REGISTERED PATENT AGENT]

Capacity: [APPLICANT / AUTHORISED PATENT AGENT]

\vspace{3mm}

\textbf{7. ABSTRACT OF THE INVENTION}

The abstract is provided on a separate page of this
complete specification.

\newpage

% =========================================================
% DESCRIPTION
% =========================================================

\begin{center}
    \textbf{\Large COMPLETE SPECIFICATION}\\[6pt]
    \textbf{\InventionTitle}
\end{center}

% ---------------------------------------------------------
\section{FIELD OF THE INVENTION}

Describe the technical field to which the invention
relates. Identify the relevant technical domain, system,
process, device, apparatus, or technical application.

[INSERT FIELD OF INVENTION]

% ---------------------------------------------------------
\section{BACKGROUND OF THE INVENTION}

Describe the existing technical approaches relevant to
the invention, including known systems, methods, devices,
processes, or technologies.

Explain the technical limitations of the existing
approaches. Support specific statements about prior art
with appropriate references wherever available.

[INSERT BACKGROUND AND PRIOR-ART DISCUSSION]

% ---------------------------------------------------------
\section{PROBLEMS ADDRESSED BY THE INVENTION}

Identify the technical problem or problems addressed by
the invention.

Explain why the existing approaches do not adequately
address these problems and identify the technical
requirements that motivate the invention.

[INSERT TECHNICAL PROBLEMS]

% ---------------------------------------------------------
\section{OBJECTS OF THE INVENTION}

The principal object of the invention is to
[INSERT PRIMARY OBJECTIVE].

Further objects of the invention include:

\begin{enumerate}
    \item To provide [INSERT OBJECTIVE].
    \item To improve [INSERT TECHNICAL CAPABILITY].
    \item To reduce [INSERT TECHNICAL LIMITATION].
    \item To enable [INSERT TECHNICAL FUNCTION].
\end{enumerate}

Include only objectives supported by the actual invention.

% ---------------------------------------------------------
\section{SUMMARY OF THE INVENTION}

Describe the invention at a level sufficient to explain
its essential technical features and how they address
the identified problem.

The invention provides [INSERT OVERVIEW OF THE INVENTION].

In one embodiment, the invention comprises
[INSERT ESSENTIAL COMPONENTS OR PROCESS STEPS].

In another embodiment, [INSERT OPTIONAL OR ALTERNATIVE
IMPLEMENTATION, IF APPLICABLE].

Explain the technical interaction between the components
or steps and the resulting technical effect.

The summary must be consistent with the detailed
description and the claims.

% ---------------------------------------------------------
\section{BRIEF DESCRIPTION OF THE DRAWINGS}

Where drawings are included, identify each figure briefly.

\begin{description}[style=nextline]
    \item[FIG. 1] [Overall architecture or system overview.]
    \item[FIG. 2] [Detailed block diagram or component arrangement.]
    \item[FIG. 3] [Flowchart of the method.]
    \item[FIG. 4] [Illustration of an embodiment or operational sequence.]
\end{description}

Delete entries for figures that are not included.

% ---------------------------------------------------------
\section{DEFINITIONS AND TERMINOLOGY}

Define specialised terms, symbols, abbreviations, and
technical expressions used in the specification where
their meanings may otherwise be unclear.

\begin{description}[style=nextline]
    \item[Term 1:] [Definition.]
    \item[Term 2:] [Definition.]
    \item[Abbreviation:] [Expanded form and meaning.]
\end{description}

% ---------------------------------------------------------
\section{DETAILED DESCRIPTION OF THE INVENTION}

Describe the invention fully and particularly, including
its components, operation, use, and the manner in which
it is performed.

\subsection{Overall Architecture}

Describe the overall arrangement of the invention,
including its constituent components, modules, interfaces,
and their relationships.

[INSERT ARCHITECTURAL DESCRIPTION]

\subsection{Components or Modules}

\begin{enumerate}
    \item \textbf{First component:}
    [Describe its structure, inputs, outputs, and function.]

    \item \textbf{Second component:}
    [Describe its structure, inputs, outputs, and function.]

    \item \textbf{Third component:}
    [Describe its structure, inputs, outputs, and function.]
\end{enumerate}

\subsection{Operational Method}

Describe the operation of the invention in a logical
sequence. Identify the relevant components and data
exchanged at each stage.

\begin{enumerate}
    \item \textbf{Step S101:}
    [First operation or acquisition step.]

    \item \textbf{Step S102:}
    [Processing, transformation, or intermediate step.]

    \item \textbf{Step S103:}
    [Core technical operation.]

    \item \textbf{Step S104:}
    [Output generation, control, or implementation step.]
\end{enumerate}

Adapt the steps to the actual invention. Do not imply that
every embodiment must perform optional steps.

\subsection{Mathematical or Algorithmic Formulation}

Where applicable, provide the mathematical formulation
and explain each variable, parameter, constraint, and
computational operation.

For example:

\begin{equation}
    y = f(x;\theta)
\end{equation}

where $x$ represents [INPUT], $\theta$ represents
[PARAMETERS], and $y$ represents [OUTPUT].

[INSERT THE ACTUAL FORMULATION, IF APPLICABLE]

Explain how the formulation participates in the claimed
technical solution. Mathematical expressions alone do
not replace a complete description of the invention.

\subsection{Embodiments and Alternatives}

Describe practical embodiments and technically
supported alternatives, including variations in
components, parameters, configurations, or operating
conditions.

[INSERT EMBODIMENTS]

% ---------------------------------------------------------
\section{BEST METHOD OF PERFORMING THE INVENTION}

Disclose the best method of performing the invention
known to the applicant and for which protection is sought.

Provide sufficient implementation detail, including
relevant component configurations, operating parameters,
process conditions, and dependencies.

[INSERT BEST METHOD KNOWN TO THE APPLICANT]

% ---------------------------------------------------------
\section{WORKING EXAMPLES AND EXPERIMENTAL VALIDATION}

Where applicable, provide reproducible working examples
demonstrating how the invention operates.

\subsection{Experimental Setup}

[Describe the apparatus, computing environment,
materials, data, or test conditions.]

\subsection{Procedure}

[Describe the actual procedure followed.]

\subsection{Observed Results}

[Report measured results and the conditions under which
they were obtained. Distinguish actual observations from
anticipated results.]

\begin{table}[H]
\centering
\caption{Illustrative structure for reporting experimental results}
\begin{tabular}{|p{45mm}|p{35mm}|p{45mm}|}
\hline
\textbf{Test or condition} &
\textbf{Measured parameter} &
\textbf{Observed result} \\
\hline
[Condition 1] & [Metric] & [Actual result] \\
\hline
[Condition 2] & [Metric] & [Actual result] \\
\hline
\end{tabular}
\end{table}

Delete the table if it is not relevant. Do not insert
fabricated results or present illustrative values as
experimental findings.

% ---------------------------------------------------------
\section{INDUSTRIAL APPLICABILITY}

Explain how the invention can be made or used in an
industry, technical process, or practical application,
as applicable.

Identify the relevant application and explain how the
invention performs the stated technical function.

[INSERT INDUSTRIAL APPLICABILITY]

% ---------------------------------------------------------
\section{ADVANTAGES AND TECHNICAL EFFECTS}

Describe the technical advantages supported by the
disclosure and, where available, experimental evidence.

Examples of possible categories include:

\begin{enumerate}
    \item Improvement in [TECHNICAL PERFORMANCE].
    \item Reduction in [RESOURCE CONSUMPTION OR LATENCY].
    \item Improvement in [RELIABILITY OR ROBUSTNESS].
    \item Simplification of [TECHNICAL PROCESS].
\end{enumerate}

Retain only advantages that can be substantiated.

% =========================================================
% CLAIMS
% =========================================================

\newpage

\begin{center}
    \textbf{\Large CLAIMS}
\end{center}

I/We claim:

\begin{enumerate}[label=\arabic*.]
    \item A [SYSTEM / DEVICE / METHOD / PROCESS] for
    [TECHNICAL PURPOSE], comprising:

    \begin{enumerate}[label=(\alph*)]
        \item [FIRST ESSENTIAL FEATURE];
        \item [SECOND ESSENTIAL FEATURE]; and
        \item [THIRD ESSENTIAL FEATURE],

        wherein [DEFINE THE TECHNICAL RELATIONSHIP OR
        OPERATION THAT CHARACTERISES THE INVENTION].
    \end{enumerate}

    \item The [SYSTEM / DEVICE / METHOD / PROCESS] as
    claimed in claim 1, wherein [INSERT A SPECIFIC
    ADDITIONAL FEATURE].

    \item The [SYSTEM / DEVICE / METHOD / PROCESS] as
    claimed in claim 1, wherein [INSERT ANOTHER
    SUPPORTED FEATURE].

    \item The [SYSTEM / DEVICE / METHOD / PROCESS] as
    claimed in claim 1, wherein [INSERT A SUPPORTED
    CONFIGURATION OR OPERATING CONDITION].

    \item A method for [TECHNICAL PURPOSE], comprising:

    \begin{enumerate}[label=(\alph*)]
        \item [FIRST METHOD STEP];
        \item [SECOND METHOD STEP]; and
        \item [THIRD METHOD STEP].
    \end{enumerate}

    wherein [STATE THE RELEVANT TECHNICAL RELATIONSHIP,
    IF APPLICABLE].
\end{enumerate}

% IMPORTANT:
% These are structural examples, not proposed legal claims.
% Adapt claim categories and dependencies to the invention.
% Ensure every claim is clear, succinct, fairly based on
% the description, and supported by the disclosure.
% Check unity of invention and statutory patentability.
% Do not retain claims for subject matter not disclosed.

% =========================================================
% SIGNATURE
% =========================================================

\newpage

\textbf{DATE AND SIGNATURE}

Date: \FilingDate

Place: [PLACE]

Signature: \rule{55mm}{0.4pt}

Name: [APPLICANT OR AUTHORISED REGISTERED PATENT AGENT]

Capacity: [APPLICANT / AUTHORISED PATENT AGENT]

% =========================================================
% DRAWINGS
% =========================================================

\newpage

\begin{center}
    \textbf{\Large DRAWINGS}
\end{center}

% Replace the following placeholders with appropriately
% prepared patent drawings. Remove this section if no
% drawings are required, subject to applicable requirements.

\begin{figure}[H]
    \centering
    % \includegraphics[width=0.85\textwidth]{figure1}
    \fbox{\parbox[c][45mm][c]{0.8\textwidth}{
        \centering INSERT FIGURE 1 HERE
    }}
    \caption*{FIG. 1}
\end{figure}

\begin{figure}[H]
    \centering
    % \includegraphics[width=0.85\textwidth]{figure2}
    \fbox{\parbox[c][45mm][c]{0.8\textwidth}{
        \centering INSERT FIGURE 2 HERE
    }}
    \caption*{FIG. 2}
\end{figure}

% Ensure drawings and reference numerals are consistent
% with the description and claims. Follow current Rule 15
% requirements for drawings and avoid decorative graphics.

% =========================================================
% ABSTRACT
% =========================================================

\newpage

\begin{center}
    \textbf{\Large ABSTRACT}
\end{center}

\textbf{\InventionTitle}

[Provide a concise technical summary of the invention,
normally not exceeding 150 words. Identify the technical
field, the core technical solution, its principal
technical advancement, and its principal use. Include
relevant reference numerals if drawings are used.]

\vspace{4mm}

\textbf{Figure to accompany the abstract, if any:}
[FIG. NUMBER]

% =========================================================
% OPTIONAL DISCLOSURES / ANNEXURES
% Include only when relevant and legally required.
% =========================================================

\newpage

\section{ADDITIONAL DISCLOSURES, IF APPLICABLE}

\subsection{Biological Material}

If applicable, provide the required identifying details
of biological material, depositary institution, deposit
number and date, and source and geographical origin,
as required by law.

[INSERT DETAILS OR DELETE THIS SUBSECTION]

\subsection{Sequence Listing}

If applicable, provide a sequence listing in the format
and manner required by the current rules.

[INSERT DETAILS OR DELETE THIS SUBSECTION]

\subsection{Other Required Technical Information}

Provide any other technical information or disclosures
required for the particular invention and applicable
Indian patent requirements.

[INSERT DETAILS OR DELETE THIS SUBSECTION]

\end{document}
